\section{Introduction}
A service agreement may be renewed under several customer conditions and then handed to a common execution system. The renewal desk knows the customer's accepted obligation. The execution system may receive only a short service code. This separation makes history a design resource: retaining fewer codes simplifies deployment, but forces some accepted continuations to share an action. The relevant question is not merely how to represent an unconstrained optimum. It is how much service value each additional code buys, and how that value depends on when participation is checked.

We study that question in a renewal model with heterogeneous continuation caps and a common terminal operating state. The provider commits to deliver an expected total equal to the weighted mean of the branch caps. This saturated promise is substantive: it makes every branch's accepted continuation bind. Intermediate service is costly, and terminal service has an increasing concave operating return. A writable code is the only branch-specific information retained at the common terminal state. The public state, read-only program, and numerical precision are distinct resources, not additional uncharged history channels.

Two institutions yield different design prescriptions. Under expected participation, the customer approves the continuation before a private service-code draw. Adjacent-code lotteries can improve value even when a deterministic codebook has already been optimized. Under realization-by-realization participation, every draw must respect the same cap. Randomization then provides no improvement at any code budget in this renewal architecture. Thus the value of randomization is an institutional choice, not a universal property of limited-memory control. A concrete advance-service and dispatch-ticket mechanism states who observes the draw, when exit is available, and which constraints a realization may exceed. A bounded-overrun variant connects the two extremes without interpreting an expected cap as a hard billing limit.

Our central optimization result jointly chooses the service levels and their lotteries. For quadratic terminal returns and heterogeneous quadratic intermediate costs, an optimal codebook can anchor every nonhighest level at a branch cap. The highest level remains continuous. We give an exact example in which restricting it to the caps misses the global optimum. This structural reduction yields a finite dynamic program with exact optimization of each continuous terminal interval. The result is a global frontier for every requested memory budget, rather than a fixed-codebook evaluation or a stationary-point calculation.

The computational consequence rests on a second model-specific fact. After optimizing the highest level continuously within an interval, the terminal-cost array still satisfies a Monge inequality. The proof uses a single-crossing identity in which heterogeneous intermediate curvatures cancel; convexity of the individual minimizations alone would not establish the result. The resulting one-sided feasible domains are classical easy-direction staircase matrices. We make that identification explicit and apply standard matrix searching, with complete tie and reachability proofs matched to the implementation. The linear-work guarantee is in exact rational arithmetic on already ordered inputs, not a new general staircase-search algorithm or a claim of uniform wall-time superiority.

The complete budget frontiers become primitives of a resource-allocation decision. A designer can price writable bits, a chosen read-only representation, evaluation work, and amortized compilation. This accounting prevents a small writable alphabet from being mistaken for a small total implementation. The electronic companion retains the broader compact-response compiler, tight response-size bounds, piecewise-quadratic closure, behavioral minimization, and shared-table comparison under their original assumptions. They explain where exact memory requirements come from; they are not additional headline results competing with the service-design frontier in the present article.

We evaluate the remaining algorithmic claims separately. Exact regression tests check costs, reconstructed controllers, and the all-submatrices completion condition. Substituting an independently authored matrix-search routine isolates implementation dependence. A standard linear-programming solver checks finite reduced network formulations. Larger same-input experiments measure runtime and memory rather than infer practical speed from an asymptotic bound. These experiments are synthetic verification and performance studies, not calibrated evidence about customer demand or willingness to accept a lottery.
